\chapter{As leis de Newton}\label{ch:g12:newtons-laws}

Contrabandeie uma balança de banheiro para dentro de um elevador. Quando a cabine
arranca para cima, o ponteiro marca três \omterm{def:g10:orders-of-magnitude:unit}{quilogramas} a mais; perto do último andar,
marca a menos; entre as duas coisas, a pura verdade --- a bordo nada mudou além do
movimento. Este capítulo enuncia as três leis por trás desse ponteiro e as forja no
método que governa toda a mecânica: escolha um sistema, desenhe as \omterm{def:g10:inertia:force}{forças} dele,
projete $\sum \vect F = m\vect a$.

\section{Quantidade de movimento e a primeira lei tornada exata}

\begin{definition}[Quantidade de movimento]\label{def:g12:newtons-laws:momentum}
A \emph{quantidade de movimento}\index{quantidade de movimento} de um corpo de massa
$m$ (\unit{kg}) que se move com velocidade $\vect v$ é o vetor
\[
\vect p = m \vect v \qquad (\unit{kg.m/s}):
\]
a velocidade, ponderada pela massa; para um sistema, a soma vetorial sobre as partes
dele.
\end{definition}

\begin{definition}[Referencial inercial]\label{def:g12:newtons-laws:inertial}
Um \emph{referencial inercial}\index{referencial inercial} é aquele em que todo corpo
isolado --- sem \omterm{def:g10:inertia:force}{forças}, ou com \omterm{def:g10:inertia:compensate}{forças que se compensam} --- se move em linha reta com
velocidade constante (\cref{ch:g10:inertia}). O solo é um deles com excelente
precisão, assim como qualquer cabine em \omterm{def:g12:kinematics-2d:uniform}{movimento retilíneo uniforme} sobre ele; um
trem freando ou um carrossel girando não são: ali, a bagagem solta acelera sem que
ninguém a empurre.
\end{definition}

\begin{theorem}[Primeira lei de Newton]\label{thm:g12:newtons-laws:first}\index{primeira lei de Newton}
Num \omterm{def:g12:newtons-laws:inertial}{referencial inercial}, a \omterm{def:g12:newtons-laws:momentum}{quantidade de movimento} de um sistema isolado não muda:
$\vect p$ é constante, ou seja, $d\vect p/dt = \vect 0$.
\end{theorem}
\admitted

\begin{remark}[A inércia, afiada duas vezes]\label{rem:g12:newtons-laws:sharpened}
É o \omterm{thm:g10:inertia:principle}{princípio da inércia} do \cref{ch:g10:inertia}, promovido duas vezes: ele diz
\emph{onde} vale --- define os referenciais inerciais como aqueles em que vale --- e
protege um vetor, $m\vect v$, cuja contabilidade rende no fim do capítulo.
\end{remark}

\section{A segunda lei de Newton}

\begin{theorem}[Segunda lei de Newton]\label{thm:g12:newtons-laws:second}\index{segunda lei de Newton}
Num \omterm{def:g12:newtons-laws:inertial}{referencial inercial}, a \omterm{def:g11:forces-and-motion:netforce}{força resultante} sobre um corpo é igual à taxa de
variação da sua \omterm{def:g12:newtons-laws:momentum}{quantidade de movimento}; para massa constante, como
$\vect p = m\vect v$,
\[
\sum \vect F = \frac{d\vect p}{dt} = m\,\frac{d\vect v}{dt} = m \vect a .
\]
\end{theorem}
\admitted

\begin{remark}[Uma lei da natureza e a unidade dela]\label{rem:g12:newtons-laws:nature}
Não há demonstração: esta é uma lei da natureza, pesada contra a experiência por três
séculos. O ano passado (\cref{ch:g11:forces-and-motion}) deu a sombra dela,
$\Delta v \propto F\,\Delta t/m$; a derivada (\cref{ch:g12:kinematics-2d}) a torna
exata a cada instante --- as letras miúdas esperam no volume do primeiro ano de
graduação. A \omterm{def:g10:inertia:force}{força} ganha uma \omterm{def:g10:orders-of-magnitude:unit}{unidade}, $\unit{N} = \unit{kg.m/s^2}$; e um corpo sujeito
apenas ao seu \omterm{def:g10:universal-gravitation:weight}{peso} obedece a $m\vect a = m\vect g$: toda massa cai com
$\vect a = \vect g$, e $g = \qty{9.81}{N/kg} = \qty{9.81}{m/s^2}$ --- as duas leituras
são uma só.
\end{remark}

\begin{example}[Força média, sem precisar de detalhes]\label{ex:g12:newtons-laws:average}
Um carro de \qty{1300}{kg} vai do repouso a \qty{100}{km/h} (\qty{27.8}{m/s}) em
\qty{8.0}{s}: seja qual for a história instante a instante, a \omterm{def:g11:forces-and-motion:netforce}{força resultante} média é
$\Delta p/\Delta t = 1300 \times 27.8/8.0 \approx \qty{4.5e3}{N}$ --- a \omterm{def:g12:newtons-laws:momentum}{quantidade de
movimento} transforma uma pergunta sobre \omterm{def:g10:inertia:force}{forças} numa contabilidade de antes e depois.
\end{example}

\section{A terceira lei de Newton}

\begin{theorem}[Terceira lei de Newton]\label{thm:g12:newtons-laws:third}\index{terceira lei de Newton}
Se um corpo $A$ exerce uma \omterm{def:g10:inertia:force}{força} $\vect F_{A \to B}$ sobre um corpo $B$, então $B$
exerce sobre $A$ a \omterm{def:g10:inertia:force}{força} $\vect F_{B \to A} = -\vect F_{A \to B}$: mesma reta de ação,
mesmo módulo, sentidos opostos --- a cada instante, qualquer que seja o movimento,
por contato ou a distância.
\end{theorem}
\admitted

\begin{example}[O livro, a mesa e a Terra]\label{ex:g12:newtons-laws:book}
Um livro repousa sobre uma mesa: duas interações, dois pares. Gravitacional: a Terra
puxa o livro para baixo ($\vect P$), o livro puxa a Terra para cima. De contato: a
mesa empurra o livro para cima ($\vect N$), o livro pressiona a mesa para baixo.
Caminhar é o mesmo negócio: empurre o chão para trás; a reação dele é a única \omterm{def:g10:inertia:force}{força}
que impulsiona você para a frente.
\end{example}

\begin{omfigure}
\begin{tikzpicture}[scale=0.9, font=\small]
  \draw[thick] (0,0) -- (4.6,0); \draw (0.5,0) -- (0.5,-1.4) (4.1,0) -- (4.1,-1.4);
  \draw[fill=omExo!15] (1.6,0) rectangle (2.6,0.6);
  \draw[-Stealth, omThm, very thick] (1.9,0.3) -- ++(0,-1.7) node[below] {$\vect P$};
  \draw[-Stealth, omProp, very thick] (2.3,0.3) -- ++(0,1.7) node[above] {$\vect N$};
  \draw[-Stealth, omProp, very thick] (3.4,-0.05) -- ++(0,-1.7) node[below] {$-\vect N$};
  \draw[thick] (0,-3.1) -- (4.6,-3.1);
  \foreach \x in {0.4,0.8,...,4.5} \draw[gray] (\x,-3.1) -- (\x-0.25,-3.4);
  \node[above right] at (3.9,-3.05) {Terra};
  \draw[-Stealth, omThm, very thick] (0.9,-3.1) -- ++(0,1.7) node[above] {$-\vect P$};
\end{tikzpicture}

{\small Um livro, duas interações, quatro flechas: o par gravitacional (a Terra puxa o
livro, o livro puxa a Terra) e o par de contato (a mesa empurra o livro, o livro
pressiona a mesa). $\vect P$ e $\vect N$ \emph{não} são parceiros: eles agem sobre o
mesmo livro --- e se cancelam apenas porque a aceleração dele é nula.}
\end{omfigure}

\section{O método: escolher, desenhar, projetar}

\begin{method}[Resolver um problema de dinâmica]\label{met:g12:newtons-laws:method}
\begin{enumerate}
  \item \emph{Sistema}: nomeie o corpo (ou o conjunto de corpos) estudado.
  \item \emph{\omterm{def:g10:relative-motion:frame}{Referencial}}: escolha um \omterm{def:g12:newtons-laws:inertial}{referencial inercial} e eixos adequados ao
        movimento --- eixos da rampa, ou a \omterm{def:g12:kinematics-2d:frenet}{base de Frenet} para uma curva
        (\cref{ch:g12:kinematics-2d}).
  \item \emph{Inventário}: liste todas as \omterm{def:g10:inertia:force}{forças} --- o \omterm{def:g10:universal-gravitation:weight}{peso} e depois uma \omterm{def:g10:inertia:force}{força} por
        contato, nada além disso --- e desenhe o diagrama de corpo livre: um ponto,
        uma flecha por \omterm{def:g10:inertia:force}{força}.
  \item \emph{Projete}: escreva $\sum \vect F = m\vect a$ e projete em cada eixo.
  \item \emph{Resolva e confira}: \omterm{def:g10:orders-of-magnitude:unit}{unidades}, sinais, casos-limite ($\mu \to 0$,
        $\alpha \to 0$, \dots).
\end{enumerate}
\end{method}

\begin{definition}[Força normal e atrito]\label{def:g12:newtons-laws:friction}
Uma superfície age sobre um corpo por meio de duas componentes: a \emph{força
normal}\index{força normal} $\vect N$, perpendicular à superfície, e o \omterm{def:g10:inertia:inventory}{atrito}
$\vect f$, tangencial. A experiência dá as leis do \omterm{def:g10:inertia:inventory}{atrito}: enquanto o corpo não
desliza, o \omterm{def:g11:forces-and-motion:friction}{atrito estático} se ajusta até um teto, $f_s \leq \mu_s N$; uma vez que ele
desliza, o \omterm{def:g11:forces-and-motion:friction}{atrito cinético} vale $f_k = \mu_k N$, contra o deslizamento. Os
\emph{coeficientes de atrito}\index{coeficiente de atrito} $\mu_s$ e $\mu_k$
($\mu_k \leq \mu_s$) dependem apenas dos dois materiais.
\end{definition}

\begin{example}[A rampa, duas vezes]\label{ex:g12:newtons-laws:incline}
Um bloco desliza por uma rampa de ângulo $\alpha$; eixos: $x$ rampa abaixo, $y$
perpendicular. Em $y$: $N = mg\cos\alpha$. Em $x$, sem \omterm{def:g10:inertia:inventory}{atrito}:
$ma = mg\sin\alpha$, logo $a = g\sin\alpha$ --- sem a massa; em $\alpha = 30^\circ$,
$a = \qty{4.9}{m/s^2}$. Com \omterm{def:g11:forces-and-motion:friction}{atrito cinético}:
$a = g(\sin\alpha - \mu_k\cos\alpha)$; para $\mu_k = 0.20$,
$a = 9.81\,(0.500 - 0.173) = \qty{3.2}{m/s^2}$. E ele só ficou parado antes disso
enquanto o \omterm{def:g11:forces-and-motion:friction}{atrito estático} necessário coube sob o seu teto:
$\tan\alpha \leq \mu_s$.
\end{example}

\begin{omfigure}
\begin{tikzpicture}[scale=1.0, font=\small]
  \draw[gray!60] (-0.3,0) -- (6.4,0); \draw[thick] (0,0) -- (30:6.2);
  \draw (1.6,0) arc (0:30:1.6); \node at (14:2.0) {$\alpha$};
  \begin{scope}[shift={(30:3.6)}, rotate=30]
    \draw[fill=omExo!15] (-0.6,0) rectangle (0.6,0.75);
    \coordinate (G) at (0,0.375); \fill (G) circle (1.5pt);
    \draw[-Stealth, omDef, dashed, thick] (G) -- ++(-2.3,0) node[above] {$x$};
    \draw[-Stealth, omDef, dashed, thick] (G) -- ++(0,1.9) node[right] {$y$};
    \draw[-Stealth, omThm, very thick] (G) -- ++(0,1.5) node[left] {$\vect N$};
    \draw[-Stealth, omThm, very thick] (G) -- ++(1.35,0) node[above] {$\vect f$};
  \end{scope}
  \draw[-Stealth, omThm, very thick] (G) -- ++(0,-1.9) node[below] {$\vect P$};
  \draw[-Stealth, omExo, dashed] (G) -- ++(210:0.95) (G) -- ++(-60:1.645);
\end{tikzpicture}

{\small Diagrama de corpo livre na rampa: \omterm{def:g10:universal-gravitation:weight}{peso} $\vect P$, \omterm{def:g12:newtons-laws:friction}{força normal} $\vect N$ e
\omterm{def:g10:inertia:inventory}{atrito} $\vect f$ rampa acima, para um bloco que desliza para baixo. Os eixos ao longo
da rampa e perpendicular a ela decompõem $\vect P$ em $mg\sin\alpha$ e
$mg\cos\alpha$ (tracejados).}
\end{omfigure}

\begin{definition}[Peso aparente]\label{def:g12:newtons-laws:apparent}
O \emph{peso aparente}\index{peso aparente} de um corpo é a \omterm{def:g12:newtons-laws:friction}{força normal} entre ele e o
seu apoio --- o que uma balança sob ele marca.
\end{definition}

\begin{example}[O elevador]\label{ex:g12:newtons-laws:elevator}
Um passageiro de \qty{70}{kg} está sobre uma balança num elevador de aceleração
vertical $a_z$ (positiva para cima): $N - mg = ma_z$, logo $N = m(g + a_z)$. Puxando
para cima a $a_z = \qty{1.5}{m/s^2}$: \qty{792}{N}, mais pesado; freando perto do
topo: \qty{582}{N}, mais leve; cruzeiro estável: $mg = \qty{687}{N}$ --- a primeira
lei; em queda livre, $a_z = -g$: $N = 0$, ausência de \omterm{def:g10:universal-gravitation:weight}{peso} sobre uma balança.
\end{example}

\begin{example}[Dois blocos e uma roldana]\label{ex:g12:newtons-laws:pulley}
Um bloco $m_1 = \qty{4.0}{kg}$ sobre uma mesa sem \omterm{def:g10:inertia:inventory}{atrito} é amarrado por uma corda
leve, passando por uma roldana ideal, a um bloco pendurado $m_2 = \qty{1.0}{kg}$. A
corda transmite a mesma tensão $T$ nas duas pontas; os blocos compartilham um único
módulo de aceleração $a$. Uma segunda lei para cada um, ao longo do respectivo
movimento: $m_1 a = T$ e $m_2 a = m_2 g - T$, logo
$a = m_2\,g/(m_1 + m_2) = \qty{1.96}{m/s^2}$ e $T = m_1 a = \qty{7.8}{N}$ --- menor
do que $m_2 g = \qty{9.8}{N}$, como tem de ser para que $m_2$ acelere para baixo.
\end{example}

\begin{omfigure}
\begin{tikzpicture}[scale=0.95, font=\small]
  \draw[thick] (0,1.1) -- (4.0,1.1); \draw (0.4,1.1) -- (0.4,-0.4) (3.6,1.1) -- (3.6,-0.4);
  \draw[fill=omExo!15] (1.0,1.1) rectangle (2.0,1.75); \node at (1.5,1.42) {$m_1$};
  \draw (4.0,1.1) -- (4.25,1.35);
  \draw[fill=white] (4.25,1.35) circle (0.25); \fill (4.25,1.35) circle (1.2pt);
  \draw[thick] (2.0,1.6) -- (4.25,1.6) (4.5,1.35) -- (4.5,0.15);
  \draw[fill=omExo!15] (4.2,-0.45) rectangle (4.8,0.15); \node at (4.5,-0.15) {$m_2$};
  \draw[-Stealth, omThm, very thick] (2.0,1.42) -- ++(1.2,0) node[pos=0.8, below] {$\vect T$};
  \draw[-Stealth, omThm, very thick] (4.5,0.15) -- ++(0,0.85) node[right] {$\vect T$};
  \draw[-Stealth, omThm, very thick] (4.5,-0.15) -- ++(0,-1.1) node[right] {$\vect P_2$};
  \draw[-Stealth, omDef, thick] (1.1,2.15) -- ++(0.8,0) node[right] {$\vect a$};
  \draw[-Stealth, omDef, thick] (5.4,0.4) -- ++(0,-0.8) node[below] {$\vect a$};
\end{tikzpicture}

{\small Dois blocos, uma corda leve, uma roldana ideal: a mesma tensão $T$ puxa as
duas pontas e os blocos compartilham um único módulo de aceleração --- uma segunda
lei para cada, duas equações, duas incógnitas.}
\end{omfigure}

\begin{example}[A curva plana]\label{ex:g12:newtons-laws:curve}
Um carro de massa $m$ faz uma curva plana de raio $R$ com \omterm{def:g12:kinematics-2d:velocity}{velocidade escalar} constante
$v$. Na \omterm{def:g12:kinematics-2d:frenet}{base de Frenet} (\cref{ch:g12:kinematics-2d}) a aceleração é puramente normal,
$a_N = v^2/R$, apontada para o centro. Na vertical, $N = mg$; na horizontal, a única
\omterm{def:g10:inertia:force}{força} centrípeta disponível é o \omterm{def:g11:forces-and-motion:friction}{atrito estático} da pista sobre os pneus:
$f = mv^2/R \leq \mu_s N = \mu_s mg$, logo $v \leq \sqrt{\mu_s\, g R}$, seja qual for
a massa. Para $R = \qty{90}{m}$ em pista seca ($\mu_s = 0.70$):
$v_{\max} \approx \qty{25}{m/s}$ (\qty{89}{km/h}); no gelo o teto desaba ---
o \cref{exo:g12:newtons-laws:13} sobreleva a pista, e o
\cref{exo:g12:newtons-laws:10} faz a mesma projeção pendurada num fio.
\end{example}

\begin{omfigure}
\begin{tikzpicture}[scale=0.85, font=\small]
  \draw[gray!50, dashed] (0,0) circle (2.1); \fill (0,0) circle (1.5pt);
  \draw[dashed, omExo] (0,0) -- (150:2.1) node[midway, above] {$R$};
  \draw[fill=omExo!15, rotate around={90:(2.1,0)}] (1.6,-0.28) rectangle (2.6,0.28);
  \draw[-Stealth, omDef, thick] (2.1,0.55) -- ++(0,1.4) node[above] {$\vect v$};
  \draw[-Stealth, omThm, very thick] (2.38,0) -- ++(-1.45,0) node[pos=0.85, above] {$\vect f$};
  \node at (0.2,-2.7) {vista de cima};
\end{tikzpicture}
\qquad
\begin{tikzpicture}[scale=0.85, font=\small]
  \draw[thick] (-0.5,0) -- (3.7,0); \draw[fill=omExo!15] (0.9,0.42) rectangle (2.5,1.25);
  \draw[fill=white] (1.2,0.42) circle (0.26) (2.2,0.42) circle (0.26);
  \coordinate (C) at (1.7,0.85); \fill (C) circle (1.2pt);
  \draw[-Stealth, omThm, very thick] (C) -- ++(0,1.4) node[above] {$\vect N$};
  \draw[-Stealth, omThm, very thick] (C) -- ++(0,-1.7) node[below] {$\vect P$};
  \draw[-Stealth, omThm, very thick] (C) -- ++(-1.6,0) node[above] {$\vect f$};
  \node at (1.6,-2.7) {vista traseira (centro à esquerda)};
\end{tikzpicture}

{\small Um carro numa curva plana com \omterm{def:g12:kinematics-2d:velocity}{velocidade escalar} constante: de cima (à
esquerda) a velocidade é tangente e a \omterm{def:g11:forces-and-motion:netforce}{força resultante} aponta para o centro; de trás
(à direita) $\vect N$ equilibra $\vect P$, e o \omterm{def:g11:forces-and-motion:friction}{atrito estático} lateral da pista é toda
a \omterm{def:g10:inertia:force}{força} centrípeta.}
\end{omfigure}

\section{Sistemas isolados: a quantidade de movimento se conserva}

\begin{proposition}[Conservação da quantidade de movimento]\label{prop:g12:newtons-laws:conservation}\index{conservação da quantidade de movimento}
Para dois corpos cujas \omterm{def:g10:inertia:force}{forças} externas se anulam ou \omterm{def:g10:inertia:compensate}{se compensam}, a \omterm{def:g12:newtons-laws:momentum}{quantidade de
movimento} total $\vect p_1 + \vect p_2$ é constante, sejam quais forem as \omterm{def:g10:inertia:force}{forças} que
eles exerçam um sobre o outro.
\end{proposition}

\begin{proof}
Some as duas segundas leis:
$\frac{d}{dt}(\vect p_1 + \vect p_2) = \vect F_{2 \to 1} + \vect F_{1 \to 2} +
\vect F_{\text{ext}}$: o par interno se cancela pela terceira lei, e
$\vect F_{\text{ext}} = \vect 0$ por hipótese.
\end{proof}

\begin{example}[Recuo]\label{ex:g12:newtons-laws:recoil}
Uma patinadora de \qty{60}{kg} em repouso sobre gelo liso arremessa horizontalmente
uma bola medicinal de \qty{4.0}{kg} a \qty{6.0}{m/s}. Antes: $\vect p = \vect 0$;
depois: $0 = 60\,V + 4.0 \times 6.0$, logo $V = \qty{-0.40}{m/s}$ --- ela desliza para
trás. Nada externo a empurrou: foi a bola (terceira lei), e as contas continuam
fechando em zero.
\end{example}

\begin{remark}[Como um foguete empurra o nada]\label{rem:g12:newtons-laws:rocket}
Um foguete é uma máquina de arremessar massa. A cada segundo os motores dele lançam
uma leva de gás para trás; o gás leva embora \omterm{def:g12:newtons-laws:momentum}{quantidade de movimento} para trás, de
modo que o foguete ganha a mesma quantidade para a frente --- uma \omterm{def:g10:inertia:inventory}{propulsão} constante,
aproximadamente a massa expelida por segundo vezes a velocidade de ejeção. O foguete
empurra o próprio jato de escape, e não o ar: ele funciona melhor no espaço vazio.
\end{remark}

\section{Exercícios}

\begin{exercise}[$\star$]\label{exo:g12:newtons-laws:1}
Calcule a \omterm{def:g12:newtons-laws:momentum}{quantidade de movimento} (a) de um velocista de \qty{70}{kg} a
\qty{10}{m/s}; (b) de um carro de \qty{1300}{kg} a \qty{50}{km/h}; (c) de uma bala de
\qty{8.0}{g} a \qty{800}{m/s}. Ordene-as. Por que a mais veloz fica em último?
\end{exercise}

\begin{exercise}[$\star$]\label{exo:g12:newtons-laws:2}
Um passageiro de \qty{70}{kg} está sobre uma balança num elevador. O que ela marca
quando a cabine (a) acelera para cima a \qty{1.2}{m/s^2}; (b) sobe a constantes
\qty{2.0}{m/s}; (c) acelera para baixo a \qty{1.2}{m/s^2}? Que lei resolve o item (b)?
\end{exercise}

\begin{exercise}[$\star$]\label{exo:g12:newtons-laws:3}
Um bloco é solto numa rampa sem \omterm{def:g10:inertia:inventory}{atrito} de ângulo $30^\circ$. Calcule a aceleração dele
e depois a velocidade após deslizar \qty{3.0}{m}. Para onde foi a massa?
\end{exercise}

\begin{exercise}[$\star$]\label{exo:g12:newtons-laws:4}
Uma bola de tênis de \qty{58}{g} chega a \qty{15}{m/s} e é devolvida pela mesma reta a
\qty{25}{m/s}; o contato dura \qty{5.0}{ms}. Calcule a variação da \omterm{def:g12:newtons-laws:momentum}{quantidade de
movimento} e a \omterm{def:g10:inertia:force}{força} média das cordas; compare essa \omterm{def:g10:inertia:force}{força} com o \omterm{def:g10:universal-gravitation:weight}{peso} da bola.
\end{exercise}

\begin{exercise}[$\star$]\label{exo:g12:newtons-laws:5}
Um livro repousa sobre uma mesa. Nomeie as duas \omterm{def:g10:inertia:force}{forças} que agem sobre o livro e, para
cada uma, o par de ação e reação dela (que corpo? que sentido?). Por que o \omterm{def:g10:universal-gravitation:weight}{peso} do
livro e o empurrão da mesa não são parceiros um do outro?
\end{exercise}

\begin{exercise}[$\star\star$]\label{exo:g12:newtons-laws:6}
Um engradado de \qty{50}{kg} está num piso horizontal, com $\mu_s = 0.45$ e
$\mu_k = 0.35$. (a) Que empurrão horizontal o coloca em movimento? (b) Se esse
empurrão for mantido, qual é a aceleração dele? (c) Que empurrão o mantém em movimento
com velocidade constante?
\end{exercise}

\begin{exercise}[$\star\star$]\label{exo:g12:newtons-laws:7}
O mesmo engradado é colocado numa rampa de inclinação ajustável ($\mu_s = 0.45$,
$\mu_k = 0.35$). (a) Em que ângulo ele começa a deslizar? (b) A $30^\circ$, calcule a
aceleração dele. (c) Uma vez iniciado o movimento no ângulo do item (a), ele continua
acelerando? Por quê?
\end{exercise}

\begin{exercise}[$\star\star$]\label{exo:g12:newtons-laws:8}
Um bloco de \qty{3.0}{kg} sobre uma mesa sem \omterm{def:g10:inertia:inventory}{atrito} é amarrado, por uma roldana ideal,
a um bloco pendurado de \qty{2.0}{kg}. Calcule a aceleração e a tensão. Por que a
tensão tinha de sair menor do que o \omterm{def:g10:universal-gravitation:weight}{peso} do bloco pendurado?
\end{exercise}

\begin{exercise}[$\star\star$]\label{exo:g12:newtons-laws:9}
Uma curva de raio \qty{50}{m} é plana. Calcule a velocidade máxima de um carro para
$\mu_s = 0.70$ (pista seca) e $0.40$ (molhada), em \unit{km/h}. Por que a resposta
vale igualmente para uma motoneta e para um caminhão?
\end{exercise}

\begin{exercise}[$\star\star$]\label{exo:g12:newtons-laws:10}
Uma massa de \qty{0.30}{kg} num fio de \qty{1.2}{m} descreve um círculo horizontal,
com o fio mantendo constantes $25^\circ$ com a vertical (um pêndulo cônico). Projete
a segunda lei na vertical e depois ao longo do raio: calcule a tensão e a velocidade.
\end{exercise}

\begin{exercise}[$\star\star$]\label{exo:g12:newtons-laws:11}
Um pacote de \qty{5.0}{kg} viaja sobre uma balança num elevador de carga; as leituras:
\begin{center}\small\begin{tabular}{lccc}
fase & \qtyrange{0}{2}{s} & \qtyrange{2}{8}{s} & \qtyrange{8}{10}{s}\\
balança (\unit{N}) & 54.0 & 49.0 & 44.1
\end{tabular}\end{center}
Calcule a aceleração em cada fase e descreva um percurso possível. O elevador poderia
estar \emph{descendo} o tempo todo?
\end{exercise}

\begin{exercise}[$\star\star\star$]\label{exo:g12:newtons-laws:12}
Uma patinadora de \qty{60}{kg} em repouso sobre gelo liso arremessa horizontalmente
uma bola de \qty{4.0}{kg}. (a) A bola parte a \qty{8.0}{m/s} \emph{em relação ao
solo}: qual é a velocidade de recuo dela? (b) Os \qty{8.0}{m/s} são agora medidos
\emph{em relação a ela}: recalcule as duas velocidades. (c) Que \omterm{def:g10:inertia:force}{força} a empurrou para
trás?
\end{exercise}

\begin{exercise}[$\star\star\star$]\label{exo:g12:newtons-laws:13}
Uma curva de rodovia de raio \qty{80}{m} é projetada para \qty{20}{m/s} sem qualquer
ajuda do \omterm{def:g10:inertia:inventory}{atrito}. Mostre que a pista precisa ser sobrelevada no ângulo $\theta$ dado
por $\tan\theta = v^2/(gR)$ e calcule-o. O que o \omterm{def:g10:inertia:inventory}{atrito} precisa fazer com um carro que
faça a curva mais devagar --- ou mais rápido?
\end{exercise}

\begin{exercise}[$\star\star\star$]\label{exo:g12:newtons-laws:14}
Os motores de um foguete ejetam \qty{250}{kg} de gás por segundo a \qty{2500}{m/s}.
(a) Quanta \omterm{def:g12:newtons-laws:momentum}{quantidade de movimento} sai por segundo --- que \omterm{def:g10:inertia:inventory}{propulsão} os gases
devolvem? (b) O foguete tem \qty{5.0e4}{kg} na ignição: qual é o \omterm{def:g10:universal-gravitation:weight}{peso} dele? E a
aceleração de decolagem? (c) Por que o motor funciona igualmente bem no espaço vazio?
\end{exercise}

\begin{exercise}[$\star\star\star$]\label{exo:g12:newtons-laws:15}
Um bloco de \qty{4.0}{kg} numa mesa ($\mu_s = 0.35$, $\mu_k = 0.25$) é amarrado, por
uma roldana ideal, a uma massa pendurada $m_2$. (a) Qual é o $m_2$ mínimo que inicia o
movimento? (b) Para $m_2 = \qty{2.0}{kg}$, calcule a aceleração e a tensão.
(c) Verifique que $\mu_k = 0$ recupera a fórmula sem \omterm{def:g10:inertia:inventory}{atrito} do
\cref{exo:g12:newtons-laws:8}.
\end{exercise}

\section{Problema: o teleférico}

\begin{problem}[{Problema de fim de semana --- certificar o teleférico: a tensão que a
rampa exige em repouso, a sobretaxa da partida, o alívio sobre a torre, o preço de uma
parada de emergência --- e o único número que a placa da cabine pode prometer}]
\label{pb:g12:newtons-laws:1}
Um teleférico de montanha --- cabine mais carro, de massa vazia
$M_0 = \qty{2.0e3}{kg}$ --- rola sobre um cabo carregador inclinado a
$\alpha = 30^\circ$, puxado por um cabo tractor paralelo à rampa com tensão $\vect T$.
O \omterm{def:g10:inertia:inventory}{atrito} de rolamento é desprezível; o piso da cabine permanece horizontal o tempo
todo. A cabine tem lugar para $25$ passageiros de massa média \qty{80}{kg}; você é o
engenheiro que certifica. Salvo indicação em contrário, a cabine está cheia:
$M = \qty{4.0e3}{kg}$.

\medskip
\noindent\textbf{Parte I --- Em repouso na rampa.}
\begin{enumerate}
  \item Escolha o sistema; diga por que o \omterm{def:g10:relative-motion:frame}{referencial} da montanha serve como
        \omterm{def:g12:newtons-laws:inertial}{referencial inercial}; faça o inventário das três \omterm{def:g10:inertia:force}{forças}, com os sentidos.
  \item Que lei governa a espera? Projete-a ao longo da rampa e perpendicularmente a
        ela para exprimir $T$ e a \omterm{def:g12:newtons-laws:friction}{força normal} $N$ do cabo carregador.
  \item Calcule $T$ e $N$.
  \item Recalcule $T$ para a cabine vazia. Por que $T$ é proporcional à massa total?
  \item O cabo tractor arrebenta, com os freios soltos: mostre que a aceleração de
        descida descontrolada vale $g \sin\alpha$, com a cabine cheia ou vazia, e
        calcule-a.
\end{enumerate}

\medskip
\noindent\textbf{Parte II --- A partida.} O guincho acelera a cabine cheia a
$a = \qty{0.60}{m/s^2}$ rampa acima até ela atingir a velocidade de cruzeiro
$v = \qty{6.0}{m/s}$.
\begin{enumerate}[resume]
  \item Exprima e calcule a tensão $T'$ durante essa fase.
  \item Em que fração a partida eleva a tensão acima do valor de repouso?
  \item Quanto dura a fase, e ao longo de que distância?
  \item Uma passageira de \qty{80}{kg} está de pé no piso horizontal. Projete a
        segunda lei na vertical para exprimir e calcular o \omterm{def:g12:newtons-laws:apparent}{peso aparente} dela; compare
        com $mg$.
  \item Que \omterm{def:g10:inertia:force}{força} a acelera horizontalmente, e de que tamanho? Por que os passageiros
        se inclinam?
\end{enumerate}

\medskip
\noindent\textbf{Parte III --- Sobre a torre.} A constantes \qty{6.0}{m/s}, o cabo
carregador se curva sobre uma torre de apoio ao longo de um arco circular vertical de
raio $R = \qty{40}{m}$; no alto, a velocidade é horizontal.
\begin{enumerate}[resume]
  \item Na \omterm{def:g12:kinematics-2d:frenet}{base de Frenet} no alto (\cref{ch:g12:kinematics-2d}), dê as duas
        componentes da aceleração, com os valores delas.
  \item Projete a segunda lei ao longo da normal para exprimir e calcular a \omterm{def:g12:newtons-laws:friction}{força
        normal} $N'$ do cabo sobre o carro.
  \item A mesma pergunta para o \omterm{def:g12:newtons-laws:apparent}{peso aparente} da passageira no alto: em que fração ela
        fica mais leve?
  \item Com que velocidade $N'$ se anularia? Compare com a velocidade de cruzeiro e
        comente.
  \item Que \omterm{def:g10:inertia:force}{força} o carro exerce sobre o cabo sobre a torre --- que lei diz isso, com
        que módulo e que sentido?
\end{enumerate}

\medskip
\noindent\textbf{Parte IV --- Frenagem, um salto e a placa.}
\begin{enumerate}[resume]
  \item No trecho horizontal de chegada, uma parada de emergência leva a cabine cheia
        de \qty{6.0}{m/s} ao repouso em \qty{1.5}{s}: leia
        $\sum \vect F = d\vect p/dt$ como média ao longo da parada para calcular a
        \omterm{def:g10:inertia:force}{força} de frenagem.
  \item Que \omterm{def:g10:inertia:force}{força} horizontal precisa parar uma passageira de \qty{80}{kg}, e o \omterm{def:g11:forces-and-motion:friction}{atrito
        estático} entre sapato e piso ($\mu_s = 0.40$) consegue fornecê-la? Conclua.
  \item Na plataforma, com os freios soltos, a cabine vazia fica pendurada e imóvel; o
        último passageiro (\qty{80}{kg}) salta horizontalmente a \qty{2.5}{m/s}
        \emph{em relação à cabine}. Calcule a velocidade de recuo da cabine e a
        velocidade do passageiro em relação ao solo.
  \item Em uma frase: que dispositivo de \omterm{def:g10:inertia:inventory}{propulsão} vive desse princípio de recuo, e o
        que fixa a \omterm{def:g10:inertia:inventory}{propulsão} dele?
  \item \emph{Certificação.} O cabo tractor rompe a \qty{100}{kN}; as normas limitam a
        tensão de \omterm{def:g11:work-of-force:work}{trabalho} a um quinto disso. Mostre que o pior caso encontrado acima
        é a partida, deduza a massa total máxima e imprima a placa: quantos
        passageiros de \qty{80}{kg} podem viajar? Confira o coeficiente de segurança.
\end{enumerate}
\end{problem}
